\mychapter{Introduction}{Introduction}{}
\label{chap:introduction}

\mysection{Problem Statement}{Problem Statement}
\label{sec:problemstatement}

This project will determine whether using Dynamic Time Warping (DTW) in the frequency domain, thus 
using vector features, is able to effectively detect Bryde's Whale calls within the False Bay region.
Three feature extraction methods will be compared. These will be Short Time Fourier Transform (STFT), 
Mel-frequency Cepstral Coeffficients (MFCC) and Linear Predictive Coding (LPC).

\mysection{Objectives}{Objectives}
\label{sec:objectives}

\begin{enumerate}
    \item Implement a DTW detection system using three different feature extraction methods(STFT, MFCC and LPC). This shows whether the effectiveness 
    of the DTW system is independent from the feature extraction method.
    \item Calculate the precision, recall and F1-score for each model.
    \item Compare results and determine whether DTW is overall a good detection method and which 
    feature extraction method is best suited.
\end{enumerate}

\mysection{Scope}{Scope}
\label{sec:scope}
    What you didn't consider.

\mysection{Layout of report}{Layout of report}
\label{sec:layout}

\begin{itemize}
    \item Chapter 2 provides anu background information needed to understand how to solve the given problem as well as any helpful literature of previous 
    solutions to the problem.
    \item Chapter 3 provides an explanation of the overall setup and method used to solve the problem including what metrics were used to meet the stated objectives.
    \item Chapter 4 gives a detailed design explanation for all the functional blocks defined in the previous chapter.
    \item Chapter 5 provides evidence of results obtained after the system has been designed as well as descriptions of what the results show.
    \item Chapter 6 contains all concluding thoughts on the results and what that means for the solutions as well as improvements and recommendations for the project.
\end{itemize}


\Ac{MFT} is as good as \ac{MFT} can be, as Table \ref{tab:my_label} and Figure \ref{fig:enter-label} shows.
Have a look at my explanation in Section \ref{sec:relatedwork} on page \pageref{sec:relatedwork}, which is in Chapter \ref{chap:literature}.



\begin{table}[H]
    \caption{Your table caption}
    \footnotesize
    \centering
    \begin{tabular}{ccrp{3cm}}    % <--- define number of columns and alignment
    \toprule
         & Centered number & R-aligned number &  Width-set col\\
         & [Units] & & \\
    \midrule
        Description & 100 & 140 & \multicolumn{1}{c}{24601}\\
    \bottomrule
    \end{tabular}\label{tab:my_label}
\end{table}
